%% ma: L11-B16-0064
%% lop: 11
%% chuong: 5
%% bai: 16
%% dang: 11.16.6
%% loai: TLN
%% muc: NB
%% dap_an: "1"
%% nguon: "tai-lieu-hoc-tap-mon-toan-lop-11-theo-chuong-trinh-giao-duc-pho-thong-2018.pdf (Lớp Toán Thầy Duy, Phạm Lê Duy, THPT Chu Văn An), Chương 3 Bài 2, Phần 3, Câu 2 (tr.378)"
%% kien_thuc: "Lớp 11 Bài 16 (giới hạn của hàm số)"
%% trang_thai: cho-duyet
%% ly_do: "Dạng mới đề xuất (11.16.x), chờ giáo viên duyệt dạng; tài liệu gốc không có lời giải mẫu, đáp án do kho tự giải và kiểm bằng sympy."
\begin{ex}
Hàm Heaviside có dạng $H(t)=\begin{cases}0&\text{khi }t<0\\1&\text{khi }t\ge0\end{cases}$ thường được dùng để mô tả việc chuyển trạng thái tắt/mở của dòng điện tại thời điểm $t=0$. Tính $\lim\limits_{t\to0^-}H(t)+\lim\limits_{t\to0^+}H(t)$.
\shortans{1}
\loigiai{\sloppy Khi $t\to0^-$ thì $t<0$ nên $H(t)=0$: giới hạn $0$. Khi $t\to0^+$ thì $t>0$ nên $H(t)=1$: giới hạn $1$. Tổng bằng $1$.}
\end{ex}
